% Шаблон стандарта: «Стандарт на разработку требований к ПО», стандарты на проектирование и кодирование.
% Скопируйте в docs/0N_<имя>.tex, замените название, версию и дату, заполните разделы.
% Красные пометки \clarify{…} — места, которые надо заполнить; перед сдачей их не остаётся.
% Правила оформления — docs/00_formatting.tex. Собрать: ./build.sh docs/0N_<имя>.tex
\documentclass{tpks}
\doctitle{Стандарт на …}
\docversion{1}
\docdate{ДД.ММ.ГГГГ}

\begin{document}
\makeheader
% Лист регистрации изменений — сразу после \makeheader. Новая версия — новая строка.
\begin{changelog}
\change{1}{ДД.ММ.ГГГГ}{Первая версия}{Фамилия И. О.}{Задание к паре ДД.ММ.ГГГГ}
\end{changelog}

\section{Введение}
\subsection{Назначение}
\clarify{что регламентирует стандарт}

\subsection{Область применения}
\clarify{к каким данным и участникам относится стандарт}

\subsection{Термины и сокращения}
Определения взяты из общего глоссария проекта.
\printterms{ПО}

\subsection{Ссылки}
\begin{enumerate}
  \item \clarify{«Название документа». Что из него используется}
\end{enumerate}

\section{Правила}
% Что и как записывать: язык, поля записи, пример, который сам соблюдает правила.
\clarify{правила с примером}

\section{Описание процесса}
% Та же схема, что в планах. Окружение процесса — по слайдам Козырева.
\subsection{Окружение процесса}
\clarify{схема процесса с входами и выходами}
\subsection{Условия входа}
\clarify{}
\subsection{Входные данные}
\clarify{}
\subsection{Описание}
\clarify{}
\subsection{Условия выхода}
\clarify{}
\subsection{Выходные данные}
\clarify{}
\subsection{Роли}
\clarify{}

\section{Критерии проверки}
\clarify{свойства результата, по которым его принимают}

\end{document}
